\appendix
\section{증명}\label{app:proofs}

\subsection{정리~\ref{thm:beta-star} 완전 증명}\label{app:proof:beta-star}

증명은 네 단계로 진행되며, 가정은 (i) $P_{\mathrm{ont}} = BB^\top$는 정사영, (ii) softmax는 실해석적, (iii) baseline이 ground-truth proxy 집합에 양의 질량을 둔다는 $\pi_{\mathcal{G}} > 0$ 경계 조건뿐이다. 먼저 $\beta$-섭동된 query 아래 로짓과 softmax의 미분을 적는다. $Q_\beta := (I + \beta BB^\top) Q$로 두면 각 row는 $q_\beta = (I + \beta BB^\top) q = q + \beta BB^\top q$이고, 로짓은 $\ell_\beta(t) = q_\beta^\top k_t / \sqrt d = \ell_0(t) + \beta r_t$이다. 여기서 $r_t = \tfrac{1}{\sqrt d}\langle BB^\top q, k_t\rangle$는 $\beta$에 독립인 per-token 점수이다. 따라서 $\partial_\beta \ell_\beta(t) = r_t$. $p_\beta(t) = \mathrm{softmax}(\ell_\beta(\cdot))(t)$에 대해 softmax Jacobian $\partial p_\beta(t) / \partial \ell_\beta(s) = p_\beta(t)(\delta_{ts} - p_\beta(s))$를 체인룰로 합성하면
\begin{equation}\label{eq:softmax-grad}
\partial_\beta p_\beta(t) = p_\beta(t) \Big(\partial_\beta \ell_\beta(t) - \mathbb{E}_{p_\beta}[\partial_\beta \ell_\beta]\Big) = p_\beta(t) \big(r_t - \mathbb{E}_{p_\beta}[r]\big).
\end{equation}

$L(\beta) = \sum_{t\in\mathcal{G}} p_\beta(t)$에 \eqref{eq:softmax-grad}를 적용하고 $\beta=0$에서 평가하면 $\mathbb{E}_{p_\beta}[r]\big|_{\beta=0} = \sum_t p_0(t) r_t = \bar r$이므로
\begin{align*}
\left.\frac{dL}{d\beta}\right|_{\beta=0}
&= \sum_{t\in\mathcal{G}} p_0(t)(r_t - \bar r)
 = \sum_{t\in\mathcal{G}} p_0(t) r_t \;-\; \bar r \sum_{t\in\mathcal{G}} p_0(t).
\end{align*}
$\pi_{\mathcal{G}} = \sum_{t\in\mathcal{G}} p_0(t)$와 $\bar r_{\mathcal{G}} = \frac{1}{\pi_{\mathcal{G}}} \sum_{t\in\mathcal{G}} p_0(t) r_t$로 정리하면 $\left.\frac{dL}{d\beta}\right|_{\beta=0} = \pi_{\mathcal{G}}(\bar r_{\mathcal{G}} - \bar r)$이다. 경계 조건 $\pi_{\mathcal{G}} > 0$ 아래에서 $\operatorname{sign}\!\left(\frac{dL}{d\beta}\big|_{\beta=0}\right) = \operatorname{sign}(\bar r_{\mathcal{G}} - \bar r)$이며, 따라서 원점 근방에서 $L$을 증가시키는 국소 방향은 $\bar r_{\mathcal{G}} - \bar r$의 부호와 일치한다. $\square$

이 결과는 response = fluctuation 형태의 Callen-Welton 항등식으로 해석할 수 있다: GT 질량의 $\beta$-민감도는 baseline 분포 하에서 GT 지시함수와 온톨로지-사영 점수 $r_t$의 공분산이다. 또한 이차 미분을 같은 논법으로 계산하면 $\frac{d^2 L}{d\beta^2}\big|_{\beta=0} = \pi_{\mathcal{G}}(V_{\mathcal{G}} - \mathrm{Var}_{p_0}(r))$, 여기서 $V_{\mathcal{G}} = \frac{1}{\pi_{\mathcal{G}}}\sum_{t\in\mathcal{G}} p_0(t)(r_t - \bar r_{\mathcal{G}})^2$. 따라서 Taylor 전개가 음의 곡률을 가질 때 유한 최적 $\beta^* = -s / [\pi_{\mathcal{G}}(\mathrm{Var}_{p_0}(r) - V_{\mathcal{G}})]$이 존재하며, 두 경우 모두 $\operatorname{sign}(\beta^*) = \operatorname{sign}(s)$로 부호는 보존된다.

위 논증은 원점에서의 일차 최적성만 확립하고, 전역적으로는 $L(\beta)$의 비선형성과 softmax 포화 때문에 최적 $|\beta|$가 따로 분석되어야 한다. 본문은 이 정리를 전역 최적화 정리로 쓰지 않고, 표~\ref{tab:sign-routing}의 regime 해석 도구로만 사용한다. 또한 정리의 경험적 적용은 proxy 집합 $\mathcal{G}$의 정의에 민감하며, 그 민감도는 부록~\ref{app:G-sensitivity}에서 따로 다룬다.

\subsection{정리~\ref{thm:qk-duality}와 따름정리~\ref{cor:cache-divergence} 완전 증명}\label{app:proof:qk-duality}

정리~\ref{thm:qk-duality}는 $P_{\mathrm{ont}} = BB^\top$의 대칭성만 쓴다. 스텝 $T$에서 query $q\in\mathbb{R}^d$와 key 행렬 $K\in\mathbb{R}^{T\times d}$를 고정하고 scaled dot-product attention logit $\ell(t) = q^\top k_t/\sqrt d$를 쓴다. K-측 개입 아래 수정된 key 행은 $k'_t = k_t + \alpha P_{\mathrm{ont}}^\top k_t = k_t + \alpha P_{\mathrm{ont}} k_t$이며, 따라서 새 logit은 $\ell'_K(t) = q^\top k'_t/\sqrt d = \ell(t) + (\alpha/\sqrt d)\, q^\top P_{\mathrm{ont}} k_t$이다. Q-측 개입 아래 $q' = q + \beta P_{\mathrm{ont}} q$이고 $\ell'_Q(t) = (q')^\top k_t/\sqrt d = \ell(t) + (\beta/\sqrt d)\, q^\top P_{\mathrm{ont}}^\top k_t = \ell(t) + (\beta/\sqrt d)\, q^\top P_{\mathrm{ont}} k_t$로, 대칭성 $P_{\mathrm{ont}}^\top = P_{\mathrm{ont}}$에서 $\ell'_K(t) - \ell(t)$와 $\ell'_Q(t) - \ell(t)$는 $\alpha=\beta$일 때 모든 $t$에 대해 동일하다. softmax가 logit의 함수이므로 attention 분포도 동일하고, 같은 value 행렬 아래 출력도 동일하다. $\square$

따름정리~\ref{cor:cache-divergence}는 이 per-step 동등성이 multi-step에서는 보존되지 않음을 보인다. 스텝 $T$에서 K-측 개입이 캐시에 반영되면 캐시된 key는 $k^{\text{cache}}_t = k_t + \alpha P_{\mathrm{ont}} k_t$로 고정된다. 후속 스텝 $s > T$에서 query $q_s$는 새로 계산되지만 key는 캐시값을 재사용하므로 logit은 $\ell_s^{\text{cache}}(t) = q_s^\top k^{\text{cache}}_t/\sqrt d = q_s^\top k_t/\sqrt d + (\alpha/\sqrt d)\, q_s^\top P_{\mathrm{ont}} k_t$이다. 개입 항은 $\alpha q_s^\top P_{\mathrm{ont}} k_t / \sqrt d$로 모든 후속 스텝에 공통 인자 $\alpha$를 공유하며 누적 효과를 낸다. 반면 Q-측 개입 $q'_s = q_s + \beta P_{\mathrm{ont}} q_s$는 스텝 $s$에서 새로 적용되므로 logit 개입 항은 $\beta q_s^\top P_{\mathrm{ont}} k_t/\sqrt d$로, 이전 스텝의 개입이 현재 스텝의 key 또는 query에 남지 않는다. 따라서 스텝 수 $S$에 대해 K-측 누적 개입 크기 $\lvert\sum_{s=T}^{T+S} \alpha q_s^\top P_{\mathrm{ont}} k_t/\sqrt d\rvert$는 일반적으로 $S$에 따라 선형적으로 증가할 수 있는 반면 Q-측의 스텝별 기여는 $S$와 무관하게 per-step $|\beta|$ 스케일로 제한된다. $\square$

이 두 결과는 layer-adaptive K+Q와 signed Q-only의 관계를 operator-family 수준에서 확정한다. per-step 동등성은 ``같은 가족''이라는 주장의 수학적 근거이고, multi-step 비대칭은 왜 long-horizon에서 Q-only가 K-포함 스케줄을 앞서는지의 메커니즘이다. \S\ref{sec:experiments}의 retail 액션 수 분해 결과는 이 예측과 부합한다.

\subsection{정리~\ref{thm:bound} 증명}\label{app:proof:bound}

\emph{1단계.} softmax는 실해석적이므로 $\phi$는 smooth, Taylor의 정리로 $\phi(1)-\phi(0) = \phi'(0) + \int_0^1(1-\tau)\phi''(\tau) d\tau$.
\emph{2단계.} softmax Jacobian $\partial s_t/\partial\ell_{t'} = s_t(\delta_{tt'}-s_{t'})$로 $ds_t(\tau)/d\tau = s_t(\tau)(\alpha_t - \bar\alpha(\tau))$이고 $\phi'(0) = \sum_t s_t\alpha_t(v_t-o) = L(q,E)$ (정확 등식).
\emph{3단계.} 직접 계산으로 $\phi''(\tau) = \sum_t s_t(\tau)[(\alpha_t-\bar\alpha(\tau))^2 - \mathrm{Var}_s\alpha(\tau)] v_t$.
\emph{4단계.} 벡터 Cauchy-Schwarz로 $\|L\|^2 \le \mathrm{Var}_s\alpha(q)\cdot\mathrm{Var}_s[V](q)$, $|\alpha_t|\le Q_{\max}\rho/\sqrt d$로 $\|R\|^2 \le Q_{\max}^4 V_{\max}^2\rho^4/d^2$. $\square$

\subsection{따름정리}\label{app:corollaries}

\begin{corollary}[정지 K의 반복 편향]\label{cor:seka-repeat}
그리디 디코딩 아래 첫 단계에 $t_1\in T_{\mathrm{gt}}$가 선택되고 $t_1$이 온톨로지 부분공간에 정렬되면 정지 K가 $t_1$의 logit을 $(1+\alpha)$배로 증폭해 $\Pr[t_2=t_1\mid\text{stat}] \ge \Pr[t_2=t_1\mid\text{no\_steer}]$. $\alpha=0$일 때 등호.
\end{corollary}

\subsection{$\beta^*$ 진단의 $\mathcal{G}$ 정의 민감도}\label{app:G-sensitivity}

정리~\ref{thm:beta-star}은 $\beta=0$ 근방의 일차 부호를 $\pi_{\mathcal{G}}(\bar r_{\mathcal{G}} - \bar r)$로 환원하지만, 이 진단을 경험적으로 구현하려면 proxy 집합 $\mathcal{G}$를 고정해야 하고 그 선택에 따라 예측의 경험적 정합성이 달라진다. 본 절은 후보 세 변형을 나란히 놓고 어느 것이 현재 단계에서 검증된 진단인지, 어느 것이 열린 과제인지 투명하게 구분한다. 정리 자체의 수학적 타당성(식 \eqref{eq:softmax-grad}의 softmax gradient identity)은 $\mathcal{G}$의 정의와 무관하게 성립하며, 아래 논의는 어디까지나 이 정리를 ``배포 가능한 부호 예측기''로 쓸 수 있느냐에 대한 범위 진술이다.

첫 번째 변형은 스키마-$\mathcal{G}$로, ground-truth 도구의 JSON 스키마 스팬에 속하는 토큰 집합을 그대로 $\mathcal{G}$로 쓴다. 이 정의는 본문 §\ref{sec:theory}의 표~\ref{tab:sign-routing}이 해석적으로 사용하는 정의이지만, 직접 측정 수준에서는 $\tau^2$ Telecom의 경험적 부호(N=20 smoke test)를 재현하지 못한다. 구체적으로 Telecom에서는 모든 태스크가 음의 $s$를 산출해 일정한 ``$-$'' 예측을 내는데, 경험적 우세 방향은 $\beta = +0.05$이므로 N=20 상의 sign-agreement는 $31.2\%$(locked)에 머문다. Retail의 동일 smoke test(N=30)는 sign-agreement가 $66\%$ 수준으로 경계에 걸쳐 있으나, 이 값은 클래스 불균형의 부산물로 해석되어야 하며 스키마-$\mathcal{G}$가 Retail에서도 안정적 예측기라고 주장할 수 없다. 진단적 관측은 스키마 토큰에서 $\bar r_{\mathcal{G}}$가 $\bar r$보다 체계적으로 작게 나온다는 것인데, 이는 JSON 스키마 스팬이 prompt 전체 평균보다 온톨로지 축에 약하게 정렬되어 있기 때문이며, 따라서 단일 step 목적함수 $L(\beta) = \sum_{t\in\mathcal{G}} p_\beta(t)$가 생성 전 과정에 걸친 F1과 탈동조된다.

두 번째 변형은 logit-lens 판별적-$\mathcal{G}$이다. ground-truth 토큰을 유지하되, per-token 점수 $r_t$ 대신 정답 도구 로짓의 $\beta$-민감도 $\partial_\beta \log p(v_{c^*} \mid y_{<T})$를 그대로 써서 $s$를 정의한다. 이 변형은 $Q, K$뿐 아니라 value 경로와 unembedding까지 단일 forward pass에서 포착하므로, 생성 F1과 더 밀접하게 정렬될 것으로 기대된다. 본 논문은 이 변형을 현재 실험 사이클에서 진행 중인 falsification 테스트로 명시한다(현재 사이클, 결과 미확정). 즉 스키마-$\mathcal{G}$의 경험적 실패는 정리 자체의 반증이 아니라 proxy 선택의 문제이며, logit-lens 변형이 검증되기 전까지는 정리~\ref{thm:beta-star}의 배포-가능한 실용 부호 예측기는 ``보류''로 분류한다.

세 번째 변형은 generation-step 전 위치 집계로, 단일 step $T{-}1$이 아닌 생성되는 모든 도구 토큰 위치에서 $s$를 집계한다. 이 변형은 멀티-툴 상태 전이 과정에서 coverage 방향이 step마다 달라질 수 있다는 구조적 관찰과 부합하고, 정리~\ref{thm:no-memory}의 history-free 한계와도 자연스럽게 결합된다. 검증 여부는 logit-lens 변형과 함께 후속 작업으로 남긴다.

요약하면, 정리~\ref{thm:beta-star}는 regime 해석 도구로는 표~\ref{tab:sign-routing}에서 네 도메인 중 Airline을 제외한 세 도메인의 부호 방향과 정성적으로 일치하지만, 스키마-$\mathcal{G}$를 통한 직접 측정은 $\tau^2$ Telecom(31.2\%, locked)에서 실패하고 Retail에서도 경계이다. 본 논문은 이 한계를 투명하게 기록하고, 실용적 부호 예측기 주장은 logit-lens 변형의 검증 이후로 미룬다.

\section{MetaTool ST4 보조 지표}\label{app:multi-metric-st4}

MetaTool ST4에서는 F1 외에도 Exact가 같은 방향으로 움직인다. 본문에서 쓰는 baseline과 stronger 변형을 같은 evaluator에서 비교하면 baseline Exact은 $0.5252$이고, stronger Q-only 변형은 $0.5332$, early-K를 포함한 layer-adaptive 변형은 $0.5473$이다. 즉 ST4에서도 이득은 ``정답 도구 집합을 정확히 맞히는가''라는 더 엄격한 지표와 함께 증가한다.

\begin{table}[h]
\centering
\caption{MetaTool ST4 보조 지표. stronger Q-only와 early-K layer-adaptive 모두 Exact을 함께 끌어올린다.}
\label{tab:st4-aux}
\small
\begin{tabularx}{0.78\linewidth}{@{}Xcc@{}}
\toprule
방법 & F1 & Exact \\
\midrule
\texttt{no\_steer} & 0.7307 & 0.5252 \\
Q-only (stronger, $\beta=-0.03$) & 0.7535 & 0.5332 \\
레이어 적응형 (early-K variant) & 0.7514 & 0.5473 \\
\bottomrule
\end{tabularx}
\end{table}

\section{Signed Q 진단의 G-robustness}\label{app:asr-fallback}\label{app:G-robustness}

알고리즘~\ref{alg:asr}은 proxy ground-truth 집합 $\mathcal{G}$ 정의에 의존한다. $\mathcal{G}$를 세 가지 대안으로 정의해 robustness를 측정한다: (i) 도구 이름 단일 토큰, (ii) 도구 이름 전체 토큰 스팬의 평균, (iii) 정답 액션 스트링 전체. 세 정의 모두에서 표~\ref{tab:sign-routing}의 방향 진단이 4개 도메인 중 3개(Airline 제외)에서 유지되면 $\mathcal{G}$ 정의에 대한 민감도가 낮다고 볼 수 있다. Airline처럼 경계적인 사례는 현재 데이터만으로 부호를 강하게 해석하기보다, 본문이 채택한 layer-adaptive 기본해와 함께 읽는 것이 더 안전하다.

\section{추가 검증 축}\label{app:planned-validation}\label{app:additional-validation}

본 절은 본문 서사를 뒤집기보다 \emph{더 좁게 검증}하기 위해 남겨 둔 추가 축을 정리한다. 핵심은 세 가지다. 첫째, stationary K-family에 대한 protocol-matched head-to-head가 있어야 related-work 반론을 직접 닫을 수 있다. 둘째, Llama telecom의 positive-$\beta$ 붕괴를 failure type으로 분해하면 cross-model 논의가 더 선명해진다. 셋째, PCA-of-K 비교가 있어야 ``아무 저랭크 기저''가 아니라 ontology-aligned basis가 중요한지 더 정확히 말할 수 있다.

\begin{table}[h]
\centering
\caption{추가 검증 축과 현재 anchor. 본문 해석을 유지한 채 더 직접적인 반론을 닫기 위한 항목만 남겼다.}
\label{tab:planned-validation}
\scriptsize
\begin{tabularx}{\linewidth}{@{}lXlX@{}}
\toprule
ID & 질문 & 설정 & 현재 anchor \\
\midrule
V1 & SEKA/AdaSEKA를 동일 evaluator에서 공정 비교해 stationary K-family 반론을 직접 닫을 수 있는가 & Qwen2.5-7B, MetaTool ST4 / $\tau^2$ retail / $\tau^2$ telecom & 현재 stationary K는 각각 $-4.57$pp, $+0.58$pp, $+11.79$pp로 충분히 약하다. 다만 이 수치만으로 SEKA-family 전체를 대표한다고 말하기에는 부족하다. \\
V2 & ontology basis가 ``아무 저랭크 기저''보다 실제로 더 적합한가 & Qwen2.5-7B, PCA-of-K vs ontology, MetaTool ST4 / retail & shuffled/random control은 이미 열세이며, 남은 질문은 PCA가 ontology와 얼마나 가까운지다. \\
V3 & positive-$\beta$ Llama telecom 붕괴는 semantic error보다 format failure에 가까운가 & Llama-3.1-8B, $\tau^2$ telecom, $\beta\in\{+0.03,+0.05,+0.10\}$ verbose audit & empty 출력은 각각 145/200, 200/200, 55/200으로 이미 크다. 필요한 것은 raw generation 수준의 failure taxonomy다. \\
V4 & $L/4$ 경계는 우연한 하이퍼파라미터가 아니라 역할 분리의 최소 구현인가 & Qwen2.5-7B, layer-boundary sweep, ST4 / retail & 현재 본문은 early-K baseline 하나만 락되어 있으므로, 경계 선택의 민감도는 별도 표로 분리하는 편이 맞다. \\
V5 & 7B 결과가 model-size artifact가 아닌가 & Qwen2.5 1.5B / 3B / 14B, ST4 + telecom 핵심 셀 & 지금의 main claim은 7B와 Llama telecom 재현으로도 서지만, 추가 size sweep이 reviewer comfort를 높일 수 있다. \\
\bottomrule
\end{tabularx}
\end{table}
